\documentclass[11pt,a4paper]{article}

\usepackage[margin=2.4cm]{geometry}
\usepackage{amsmath,amssymb}
\usepackage{booktabs}
\usepackage{xcolor}
\usepackage{enumitem}
\usepackage[colorlinks=true,linkcolor=black,urlcolor=black]{hyperref}
\usepackage{parskip}
\usepackage{titlesec}
\usepackage{fancyhdr}

\definecolor{accent}{HTML}{1F4E79}
\definecolor{shade}{HTML}{F2F5F8}

\titleformat{\section}{\Large\bfseries\color{accent}}{\thesection.}{0.6em}{}
\titleformat{\subsection}{\large\bfseries}{}{0em}{}
\titlespacing*{\section}{0pt}{1.6em}{0.6em}

\pagestyle{fancy}
\fancyhf{}
\lhead{\small\color{accent}RailCast --- buffer allocation as a linear program}
\rhead{\small\thepage}
\renewcommand{\headrulewidth}{0.4pt}

\newcommand{\code}[1]{\texttt{\small #1}}

\begin{document}

\begin{center}
{\LARGE\bfseries\color{accent} Buffer allocation as a linear program}\\[0.4em]
{\large RailCast --- feasibility, formulation, and what the data supports}\\[0.8em]
{\small 8 September 2026}
\end{center}

\vspace{0.4em}
\noindent\fcolorbox{accent!30}{shade}{%
\begin{minipage}{0.97\textwidth}
\vspace{0.3em}
\textbf{Purpose.} Assess whether an existing Irish Rail delay-prediction dataset supports a
genuine mathematical-programming component: deciding how to distribute a fixed total
schedule buffer across a train's segments so that expected end-to-end lateness is
minimised. The requirement is a formulation derived from scratch --- decision variables,
objective, constraints --- and solved with a real solver, not a library call that happens
to optimise internally.
\vspace{0.3em}
\end{minipage}}

\section{What data we actually have}

Everything needed is already in the archive. \textbf{Nothing new needs collecting.}

Per stop, \code{getTrainMovementsXML} gives:

\begin{center}
\begin{tabular}{@{}lr@{}}
\toprule
\textbf{Field} & \textbf{Populated} \\
\midrule
\code{ScheduledArrival}, \code{ScheduledDeparture} & 95.4\% \\
\code{Arrival} & 69.1\% \\
\code{Departure} & 66.4\% \\
\code{AutoArrival}, \code{AutoDepart} & with every actual \\
\code{LocationOrder}, \code{LocationType} & $\approx$100\% \\
\code{TrainOrigin}, \code{TrainDestination} & $\approx$100\% \\
\bottomrule
\end{tabular}
\end{center}

\textbf{Departure times are the unlock.} They were unused until now --- the delay model only
ever touched arrivals. Having both arrival and departure lets a journey be decomposed into
\emph{running time} (departure at stop $i-1$ to arrival at stop $i$) and \emph{dwell time}
(arrival at stop $i$ to departure at stop $i$), which is exactly the decomposition buffer
allocation requires.

Measured on the train and validation window, with the journey-consistency filter applied
(583 of 16{,}776 journeys rejected):

\begin{itemize}[leftmargin=1.4em,itemsep=0.2em]
  \item \textbf{211 segments} with $\ge$30 clean observations.
  \item \textbf{240 trains} with $\ge$20 days in which \emph{every} segment is observed ---
        a complete scenario set.
  \item \textbf{Slack} (scheduled minus 5th-percentile observed running time): median
        \textbf{90\,s}, upper quartile 150\,s, maximum 642\,s. 146 of 211 segments exceed
        one minute.
  \item \textbf{Three segments have negative slack.} \code{CRLOW}$\to$\code{ATHY} is
        scheduled at 600\,s; the 5th percentile of actual runs is 702\,s. The timetable
        there is not achievable.
\end{itemize}

That last point matters: the problem is not purely academic. There is real, measurable,
unevenly distributed slack.

\subsection*{A caution about estimating the technical minimum}

A first pass used the \emph{minimum} observed running time as the technical minimum and
produced \code{ETOWN}$\to$\code{MLGAR} scheduled at 1050\,s and apparently run in 156\,s
--- roughly 580\,km/h. \code{ETOWN} is one of four stations already identified as carrying
arrival records belonging to other trains. Minimum-of-sample is not robust to that
contamination; a 5th percentile is. Any use of a raw minimum in this work will reintroduce
the same error.

\subsection*{What is not available from the four endpoints in use}

\begin{itemize}[leftmargin=1.4em,itemsep=0.15em]
  \item connections between services, or which passengers transfer where;
  \item passenger counts or loadings;
  \item true minimum technical running times (these must be estimated);
  \item headways, track occupancy, platform capacity;
  \item rolling stock or crew circulation.
\end{itemize}

\section{Which formulation the data supports}

\subsection*{Delay management: rejected}

The delay-management problem --- deciding which connections to hold and which to drop to
minimise total passenger delay --- requires two things the dataset does not hold: a
definition of which services connect, and passenger weights on each connection.

Connections could be defined heuristically (two trains at one station within $X$ minutes),
but the passenger weights would be invented, and they sit \emph{inside the objective
function}. An optimum computed against fabricated weights is not defensible. Rejected on
those grounds.

\subsection*{Buffer allocation: supported}

The data supports it well, and it is a genuine linear program.

One caveat to state plainly: this is a \textbf{counterfactual timetabling} exercise. Trains
cannot be run on a re-buffered timetable, so evaluation is necessarily \emph{simulation} ---
replaying historical primary delays through a propagation recursion under both the existing
buffer vector and the optimised one. That is a standard and defensible method, but the
recursion \emph{is} a model, and its assumptions are the honest limitations of the result
rather than a footnote.

\subsection*{A correction on scenario generation}

It is tempting to describe this as sample average approximation over the project's existing
quantile predictions. That does not hold. The model predicts \emph{total accumulated arrival
delay at a stop given a vantage stop}, not per-segment primary delay. Differencing
consecutive predictions does not recover a coherent joint scenario along a route, because
the three quantile models were never fitted to be jointly consistent across stops.

The defensible scenario source is \textbf{historical per-segment excess running time}: whole
days, so that within-day correlation is preserved by construction. That is genuine sample
average approximation.

\section{The formulation}

\subsection*{In plain terms}

A timetable contains padding --- scheduled running times are longer than the fastest a train
can actually complete the segment. That padding absorbs delay. It is currently distributed
however the timetable happens to distribute it. The question is whether, given the
\emph{same total} padding, there is a better place to put it so that trains arrive less late
on average.

The decision is how many seconds of padding to give each segment. No padding can be
invented; it can only be moved. Real days are then replayed through the result and the
lateness counted.

\subsection*{Index and data}

Stops $i = 0, 1, \dots, n$ along one train's route. Segment $i$ runs from stop $i-1$ to
stop $i$.

\[
\begin{aligned}
m_i &= \text{minimum technical running time on segment } i \\
r_i &= \text{scheduled running time on segment } i \\
b_i^0 &= r_i - m_i \\
B &= \textstyle\sum_{i=1}^{n} b_i^0 \\
\delta_i^\omega &= \max\!\left(0,\ \rho_i^\omega - m_i\right)
\end{aligned}
\]

\noindent where $\rho_i^\omega$ is the observed running time on segment $i$ on day $\omega$.
$m_i$ is estimated as the 5th percentile of observed running times, for the robustness
reason given in section 1. $b_i^0$ is the timetable's existing buffer, which is also the
baseline to beat. $B$ is the budget: the same total buffer is redistributed, never added to,
so end-to-end journey time is unchanged. $\delta_i^\omega$ is the primary delay --- the
excess running time actually incurred on that segment that day.

\subsection*{Decision variables}

\[
b_i \ge 0 \quad (i = 1 \dots n), \qquad\qquad L_i^\omega \ge 0 \quad (i = 1 \dots n,\ \omega \in \Omega)
\]

$b_i$ is the buffer allocated to segment $i$. $L_i^\omega$ is arrival lateness at stop $i$
under scenario $\omega$.

\subsection*{Objective}

\[
\min_{b,\,L}\quad \frac{1}{|\Omega|}\sum_{\omega \in \Omega}\ \sum_{i=1}^{n} w_i\, L_i^\omega
\]

\subsection*{Constraints}

\[
\begin{aligned}
L_i^\omega &\ \ge\ L_{i-1}^\omega + \delta_i^\omega - b_i && \forall i,\ \forall \omega \\
L_i^\omega &\ \ge\ 0 && \forall i,\ \forall \omega \\
\textstyle\sum_{i=1}^{n} b_i &\ \le\ B \\
0 \le b_i &\ \le\ \bar b_i && \forall i \\
L_0^\omega &\ =\ \ell_0^\omega && \text{(initial delay; } 0 \text{ if the train starts on time)}
\end{aligned}
\]

\subsection*{The linearisation, which is the part that matters}

The true dynamics are

\[
L_i^\omega = \max\!\left(0,\ L_{i-1}^\omega + \delta_i^\omega - b_i\right),
\]

that is: a train carries its delay forward, buffer eats into it, and it cannot arrive
``negatively late'' in a way that would let it depart early. That $\max$ is
piecewise-linear, not linear, and cannot enter a linear program as written.

The standard move is to replace the \emph{equality} with the two \emph{inequalities} above
and let the objective do the rest. The argument that this is exact:

\begin{enumerate}[leftmargin=1.6em,itemsep=0.25em]
  \item Each $L_i^\omega$ enters the objective with coefficient $w_i / |\Omega| > 0$, so the
        solver drives it as small as possible.
  \item Its only lower bounds are the two constraints. Pushed down, it lands exactly on
        the value $\max\bigl(0,\ L_{i-1}^\omega + \delta_i^\omega - b_i\bigr)$ that the
        true recursion specifies.
  \item $L_i^\omega$ also appears on the right-hand side of the constraint for
        $L_{i+1}^\omega$, with a \emph{positive} coefficient --- so reducing it
        \emph{relaxes} the next constraint. Both forces act in the same direction, giving
        no incentive to inflate it.
\end{enumerate}

The linear program's optimum therefore reproduces the nonlinear recursion exactly.

\medskip
\noindent\fcolorbox{accent!30}{shade}{%
\begin{minipage}{0.97\textwidth}
\vspace{0.3em}
\textbf{The failure mode to know.} This holds only while every $w_i \ge 0$. Introduce a
negative weight anywhere --- for instance to ``reward early arrival'' --- and the relaxation
breaks: the solver inflates that $L_i^\omega$ to buy an improvement elsewhere, and the
lateness variables stop representing lateness. Rewarding earliness requires a different
encoding, not a negative weight.
\vspace{0.3em}
\end{minipage}}

\subsection*{Choosing the weights}

$w_i$ is a modelling decision. Setting $w_i = 0$ for all stops but the terminus optimises
end-to-end punctuality only; $w_i = 1$ everywhere optimises for passengers alighting
anywhere along the route. Weighting by boardings would be preferable, and boardings are not
available.

\subsection*{Size and solver}

For a 29-stop DART service with 32 complete days: 28 buffer variables, $28 \times 32 = 896$
lateness variables, and roughly 1{,}820 constraints. This solves in milliseconds.

\code{scipy.optimize.linprog} with the HiGHS backend is already installed as a transitive
dependency, so no new package is required. The constraint matrix should be built sparse; its
structure is very sparse and a dense representation wastes memory for no benefit at larger
scale.

\section{Baseline and limitations}

\textbf{Baseline.} $b^0$, the timetable's own buffer vector, evaluated with the same
recursion on \textbf{held-out days}. Buffers are fitted on one set of days and evaluated on
another. Optimising and evaluating on the same days measures only the ability to fit noise.

\textbf{Limitations, to be published alongside whatever number results:}

\begin{enumerate}[leftmargin=1.6em,itemsep=0.25em]
  \item \textbf{Primary delays are assumed exogenous} --- that $\delta_i^\omega$ does not
        change when $b_i$ changes. In practice drivers pace to the timetable, so a segment
        given more buffer may simply attract a slower run. This is the largest threat to the
        result and cannot be tested with this data.
  \item \textbf{Minimum running time is estimated}, not known. The result moves with the
        choice of percentile; the sensitivity should be reported.
  \item \textbf{Single train, no network.} A re-buffered service could conflict with another
        on shared track. There are no headway or platform constraints.
  \item \textbf{Counterfactual.} The result cannot be validated against real operations, only
        simulated.
  \item \textbf{Dwell buffer is ignored} in the base formulation; buffer lives only on
        running segments.
\end{enumerate}

\textbf{One assumption that is \emph{not} required:} the program does not assume independence
between segments. Because each scenario is a whole day, within-day correlation is preserved
exactly. This is a genuine advantage of sample average approximation over a parametric
approach.

\vspace{1em}
\hrule
\vspace{0.6em}
\noindent\small\textit{If the optimiser does not beat the existing timetable, that result is
recorded and kept. A negative result honestly measured is worth more than a claim that has to
be hedged.}

\end{document}
